% ===================================================================
%  HOTO Na 15 — Kasuwa. --- Abinci.
%
%  Traduction, faite en 2026, en haoussa d'aujourd'hui, sur l'ido de
%  texto/io. Regles de la colonne : voir 10-hoto-01.tex. Une seule
%  numerotation.
%
%  LA POMME DE TERRE (61) EST « dankalin turawa », LA PATATE DES
%  EUROPEENS, ET C'EST LA TROISIEME FOIS QUE CETTE COLONNE FAIT LE
%  MEME GESTE. Le haoussa avait le dankali, la patate douce, qui
%  pousse ici depuis toujours ; l'autre tubercule est venu des Andes
%  par l'Europe, et il a pris le mot de la patate en lui ajoutant ceux
%  qui l'apportaient. Avant lui, la neige du tableau 4 etait « dusar
%  ƙanƙara », la poudre de glace, et le charbon de terre du tableau 6
%  « gawayin dutse », le charbon de bois de pierre. TROIS DERIVATIONS,
%  ZERO EMPRUNT, ET LES TROIS CHOSES SONT LES TROIS QUI SONT ARRIVEES
%  LE PLUS TARD. Quand le mot manque, cette langue ne va pas le
%  chercher : elle en prend un qu'elle a et lui ajoute d'ou vient la
%  chose ou de quoi elle est faite.
%
%  ET UNE LANGUE MUSULMANE NOMME ICI CE QUE SES LOCUTEURS NE MANGENT
%  PAS, SANS UN DETOUR. L'etal du charcutier (2 a 7) est du jambon,
%  du boudin, de l'andouille et du lard ; la colonne ecrit naman
%  alade, tsiran alade, gasasshen naman alade. Ni euphemisme, ni
%  refus, ni note. C'est la troisieme fois du livret : la giya de la
%  noce au tableau 4, la giya et le houblon au tableau 8, la
%  charcuterie ici. La langue a les mots ; c'est le manger qui est
%  interdit, pas le nommer.
%
%  ONZIEME REFUS DE LA COLONNE : LE PAIN (79). « gurasa » est le pain
%  de froment haoussa, une galette epaisse cuite au feu sur une
%  plaque, qu'on vend au marche de Kano tous les matins et qui se
%  rompt a la main. Le boulanger de la planche a un fournil voute, il
%  a petri sa pate, et il retire du four des croissants, des petits
%  pains et des miches rondes. On ecrit « burodi ».
%
%  L'ANGUILLE EST « kifin maciji », LE POISSON-SERPENT, ET LE CRABE
%  EST « ƙaguwa » TOUT COURT. La ligne est nette et elle est
%  geographique : ce qui vit dans le Niger et la Benoue a un nom
%  plein — kifi, ƙaguwa, jatan lande la crevette —, ce qui vit dans
%  une mer que la langue n'a pas frequentee se decrit par ce a quoi
%  cela ressemble. Le releve a deja vu ce partage au tableau 9 avec
%  les arbres ; il tient sur les betes de l'eau.
%
%  ET LA REGLE DES LETTRES ABSENTES A FAIT SA MEILLEURE PRISE ICI.
%  L'ido nomme au 9 une confiture « riba o quinga » — groseille ou
%  coing —, deux fruits qui n'ont pas de nom haoussa, et l'alinea
%  avait ete ecrit en laissant les mots ANGLAIS dans le texte :
%  « currant » et « quince ». Ce n'est pas une faute d'orthographe,
%  c'est un mot qu'on a oublie de traduire, et aucun autre controle
%  du releve ne l'aurait vu : renvoji.py compte des renvois, plenajo
%  compare des longueurs, html.py publie. Le q de « quince » a suffi.
%  C'est la meme espece que les mots de brouillon — « Wait »,
%  « Ordre » — laisses dans les colonnes ourdoue et tamoule, mais
%  releves ici par une regle d'ALPHABET et non par une regle de
%  ligne. On ecrit kurant et kwins, refaits comme fis, abirkot et
%  furam.
%
%  ENFIN LE MARCHE EST « kasuwa », ET IL N'Y AVAIT RIEN A PESER. La
%  kasuwa est la plus vieille institution que cette langue nomme :
%  celle de Kurmi, a Kano, tient depuis le XVe siecle, et le mot n'a
%  pas bouge. La halle couverte (1) devient « babbar rumfar kasuwa »,
%  le grand hangar du marche, parce que c'est un toit qu'on ajoute et
%  non une chose nouvelle.
% ===================================================================

%%K t15-apar-1 apar
\VUsaut{4.0mm}
\VUtitre{15.0pt}{60}{HOTO Na 15}
\VUsaut{3.0mm}

%%K t15-tit-1 sub
\VUpk{11.4pt}{40}{Kasuwa. --- Abinci.}
\VUsaut{3.0mm}

%%K t15-01-1 p
1. --- Yawancin ƴan kasuwa da suke ba mu kayayyakin da suke buƙata
domin abincinmu suna taruwa a ƙarƙashin \VUgras{babbar rumfar}
\textsuperscript{(1)} \VUgras{kasuwa mai rufi} ko a titunan da suke kusa.

%%K t15-02-1 p
2. --- \VUgras{Mai sayar da naman alade} \textsuperscript{(2)}, wanda yake
sayar da \VUgras{gasasshen naman alade} \textsuperscript{(3)}, da
\VUgras{tsiran jini} \textsuperscript{(4)}, da \VUgras{ƙaramin tsira}
\textsuperscript{(5)}, da \VUgras{tsiran tumbi} \textsuperscript{(6)} da
\VUgras{kitsen alade} \textsuperscript{(7)}, yana tsaye kusa da \VUgras{mai sayar
da man shanu da cuku} \textsuperscript{(8)}, wadda wurin baje kayanta yake
cike da \VUgras{cukun Holand} \textsuperscript{(9)} mai jan ɓawo, da
\VUgras{cukun Camembert} \textsuperscript{(10)}, da \VUgras{faifan Gruyère}
\textsuperscript{(11)} da \VUgras{gunduwoyin man shanu} \textsuperscript{(12)} sabbi.

%%K t15-03-1 p
3. --- A gabanta, \VUgras{mai sayar da ganyayyaki}
\textsuperscript{(13)} tana gabatar wa abokan cinikinta kyawawan
\VUgras{tumatir} \textsuperscript{(14)}, da \VUgras{kabejin fure}
\textsuperscript{(15)}, da \VUgras{salati} \textsuperscript{(16)}, da
\VUgras{kabeji} \textsuperscript{(17)}, da \VUgras{kabewa}
\textsuperscript{(19)}, da \VUgras{kankana} \textsuperscript{(18)} har ma da
\VUgras{naman kaza} \textsuperscript{(20)}. Amma \VUgras{ɗan aiken giya,}
wanda ya tsayar da \VUgras{motar kai kayansa} \textsuperscript{(21)} cike da
\VUgras{gangunan giya} \textsuperscript{(22)}, daidai gaban wurin baje
kayanta, yana hana abokan cinikinta kusatowa. Wata mai lambu tana
kawo mata kwando na \VUgras{artichok} \textsuperscript{(23)}.

%%K t15-tit-2 sub
\VUsaut{4.0mm}
\VUpk{10.2pt}{40}{Sayarwa da saye.}
\VUsaut{3.0mm}

%%K t15-04-1 p
4. --- A kasuwa, hargitsi yana da yawa, domin lokacin \VUgras{sayarwa
ta gwanjo} \textsuperscript{(24)} ya yi. \VUgras{Uwayen gida}
\textsuperscript{(26)}, waɗanda suke zuwa can domin sayayyarsu, suna tsayawa
suna wucewa gaban shagon \VUgras{mai sayar da tumbi}
\textsuperscript{(25)}, domin su sayi \VUgras{tumbin ciki} ko \VUgras{kan ɗan
sa.}

%%K t15-05-1 p
5. --- Wani \VUgras{mai girki} \textsuperscript{(27)} mai ƙiba yana ciniki a
kan \VUgras{kifin tuna} mai kyau ƙwarai wurin \VUgras{mai sayar da kifi}
\textsuperscript{(28)} wadda take ƙara farashinta yadda ta ga dama. Ta
karɓi yawan \VUgras{kifin maciji, da kifin lebur, da kifin kogi} da
\VUgras{karfa.} \VUgras{Kifinta na gishiri} da \VUgras{katantanwar
ruwanta} suna da sabo ƙwarai.

%%K t15-tit-3 sub
\VUsaut{4.0mm}
\VUpk{10.2pt}{40}{Kaya masu arha da masu tsada.}
\VUsaut{3.0mm}

%%K t15-06-1 p
6. --- \VUgras{Mai sayar da tsuntsaye} \textsuperscript{(29)}, wadda ta
zauna kusa da ita, tana da gaskiya ƙwarai. Idan tana sayar da
\VUgras{talotalo,} ko \VUgras{babbar agwagwa,} ko \VUgras{agwagwa} ko
\VUgras{ɗan kaza} kawai, tana neman farashi mai adalci kaɗai.

%%K t15-07-1 p
7. --- A kusurwar filin, a bene na farko, tun ba da daɗewa ba, an kafa
\VUgras{mai sayar da yadi} \textsuperscript{(30)} wanda mata masu saye kullum
suke da tabbacin samun kyakkyawan zaɓi na \VUgras{zanen tebur,} da
\VUgras{ƙananan tawul,} da \VUgras{afaro} da \VUgras{abin share kwanuka}
wurinsa. A wannan bene, \VUgras{maƙerin wuƙaƙe} \textsuperscript{(31)} ya kafa
ɗakin aikinsa. Yana wasa \VUgras{wuƙaƙen} masu sayarwa a cikin rumfar
kasuwa yana kuma yi wa masu lambu \VUgras{lauje} da \VUgras{ƙarfe} mafi
tauri.

%%K t15-tit-4 sub
\VUsaut{4.0mm}
\VUpk{10.2pt}{40}{Kayan miya.}
\VUsaut{3.0mm}

%%K t15-08-1 p
8. --- A ƙarƙashin ɗakin aikinsa, tun ba da daɗewa ba, an buɗe babban
shagon abinci. Ba \VUgras{shagon kayan miya} \textsuperscript{(32)} mai sauƙi
kuma marar muhimmanci na dā ba ne kuma, wanda yake sayar da kayayyaki
na kowa kaɗai, wato \VUgras{shayi} \textsuperscript{(33)}, da
\VUgras{cakulet} \textsuperscript{(34)}, da \VUgras{sukari gunduwa}
\textsuperscript{(37)} da \VUgras{sabulu} \textsuperscript{(38)}. A can ana sayar da
kowane iri kaya, wato \VUgras{biskit,} da \VUgras{kayan gwangwani}
\textsuperscript{(35)}, da \VUgras{giya mai zaƙi} \textsuperscript{(36)}, da
\VUgras{rum,} da \VUgras{konyak,} da \VUgras{kirsh,} da \VUgras{giyar
anisi} da \VUgras{kurasao.} A can ana samun naman daji : \VUgras{zomon
daji} \textsuperscript{(39)}, da \VUgras{tsuntsayen daji}
\textsuperscript{(40)}, da \VUgras{makwarwa} \textsuperscript{(41)}, da
\VUgras{ƙananan makwarwa} \textsuperscript{(42)}, da \VUgras{agwagwar daji}
\textsuperscript{(43)} ko \VUgras{zabo} \textsuperscript{(44)}, cikin sauƙi kamar
ƴaƴan itacen ƙasashen waje kamar \VUgras{ayaba} \textsuperscript{(46)} da
\VUgras{abarba.}

%%K t15-09-1 p
9. --- \VUgras{Ɗan aiken} \textsuperscript{(45)} mai sayar da kayan miya,
mai taimako ƙwarai, yana a shirye ya ba da abin da mutum yake so :
\VUgras{jam} \textsuperscript{(47)} na kurant ko na kwins, da \VUgras{kitse}
\textsuperscript{(48)}, da \VUgras{ƙananan gurji} \textsuperscript{(49)} ko
\VUgras{sardin} \textsuperscript{(50)} mai gishiri.

%%K t15-09-2 p
Wannan yana da amfani ƙwarai ga \VUgras{mata masu saye}
\textsuperscript{(51)}, waɗanda suke samu, kusa da kayayyaki na kowa, abin
da ya fi wuya da ƴaƴan itace da suka fi daɗi : \VUgras{fir}
\textsuperscript{(52)} mai ruwa, da \VUgras{furam} \textsuperscript{(53)}, da
\VUgras{inabi} \textsuperscript{(54)}, da \VUgras{fis} \textsuperscript{(55)}, da
\VUgras{lauje} \textsuperscript{(56)} da \VUgras{wolnat} \textsuperscript{(57)}.

%%K t15-tit-5 sub
\VUsaut{4.0mm}
\VUpk{10.2pt}{40}{Abokan ciniki.}
\VUsaut{3.0mm}

%%K t15-10-1 p
10. --- \VUgras{Shinkafa} \textsuperscript{(58)} da \VUgras{ƙananan wake}
\textsuperscript{(59)} suna kusa da akwatin \VUgras{haring}
\textsuperscript{(60)} da aka gasa da hayaƙi ; \VUgras{dankalin turawa}
\textsuperscript{(61)} da \VUgras{wake} \textsuperscript{(63)} suna kusa da
\VUgras{tuffa} \textsuperscript{(62)}, da \VUgras{ɓaure}
\textsuperscript{(64)}, da \VUgras{busasshen furam} \textsuperscript{(65)} da
\VUgras{gyaɗar Turai} \textsuperscript{(66)}. Ana samun a wannan shago sabbin
\VUgras{ƙwai} \textsuperscript{(67)} da \VUgras{zaitun} \textsuperscript{(68)} ;
shi ya sa \VUgras{kwanduna} \textsuperscript{(70)} da \VUgras{ragunan sayayya}
\textsuperscript{(73)} suke cika da sauri ƙwarai.

%%K t15-11-1 p
11. --- \VUgras{Kofi} \textsuperscript{(72)} na wannan kamfani mai kyau
ƙwarai ya sami suna mai adalci tsakanin \VUgras{ma’aikatan gida}
\textsuperscript{(69)} da \VUgras{masu girki,} domin tsoho \VUgras{mai sayar
da kayan miya} \textsuperscript{(71)} da kansa yakan soya shi da kula
ƙwarai.

%%K t15-tit-6 sub
\VUsaut{4.0mm}
\VUpk{10.2pt}{40}{Abin kallo.}
\VUsaut{3.0mm}

%%K t15-12-1 p
12. --- Ba kowa ne yake zuwa kasuwa domin ya yi sayayya ba. A cikin
zirga-zirga mai ban sha’awa a kan ƙaramar hanyar da take kai wa rumfar
kasuwa, mutum yana ganin mutane iri iri ƙwarai. Kusa da \VUgras{ɗan
aiken mahauci} \textsuperscript{(75)} mai kirki, wanda yake tura
\VUgras{amalanke} \textsuperscript{(74)} a gabansa, ga sojoji biyu masu hutu
suna alfahari ƙwarai da nuna kayan aikinsu masu ƙyalli :
\VUgras{sojan doki} \textsuperscript{(76)} da shuɗiyar \VUgras{rigarsa} da
\VUgras{hular gashin tsuntsu,} da \VUgras{sajan zuwaf,} mai gajeren wando
mai kumbura da kai a rufe da \VUgras{hular fez.}

%%K t15-13-1 p
13. --- \VUgras{Mai gasa burodi} \textsuperscript{(80)}, wanda yake tsaye
a bakin \VUgras{gidan gasa burodi} \textsuperscript{(78)}, ya gama kwaɓa
garin \VUgras{burodinsa} \textsuperscript{(79)} ya kuma ɗauko daga tanda
\VUgras{burodi mai lanƙwasa} \textsuperscript{(81)}, da \VUgras{ƙananan
burodi} \textsuperscript{(82)} da \VUgras{burodi mai zagaye}
\textsuperscript{(83)} waɗanda suke a tagar shagon.

%%K t15-14-1 p
14. --- A saman gidan gasa burodin, akwai gwaniyar \VUgras{mai ɗinkin
farar kaya} \textsuperscript{(84)} zaune, wadda take ɗaukar \VUgras{mata masu
ɗinkin ado} \textsuperscript{(85)} da yawa. Kusa da ita akwai \VUgras{mai
wanki da goge tufafi} \textsuperscript{(86)} zaune, wadda take sa sitaci a
\VUgras{farar kaya} \textsuperscript{(88)} bayan ta wanke ta kuma shanya
ta, kuma wadda take goge ta da \VUgras{ƙarfen goge tufafi}
\textsuperscript{(87)}.

%%K t15-tit-7 sub
\VUsaut{4.0mm}
\VUpk{10.2pt}{40}{Nama, kifi da ganyayyaki.}
\VUsaut{3.0mm}

%%K t15-15-1 p
15. --- A \VUgras{mahauta} \textsuperscript{(89)}, wadda take a benen ƙasa,
ana sayar da kowane iri \VUgras{nama, naman tunkiya}
\textsuperscript{(90)}, da \VUgras{naman sa} \textsuperscript{(94)} ko
\VUgras{naman ɗan sa} \textsuperscript{(92)} a matsayin \VUgras{cinyar
tunkiya, da yankan nama, da gasasshen nama} da \VUgras{yankan haƙarƙari.}
\VUgras{Ɗan aiken mahauci} \textsuperscript{(93)} yana yayyanka dukan
waɗannan nama yana kuma ware \VUgras{ƙasusuwan ɓargo}
\textsuperscript{(96)}. A ƙofar mahautar akwai \VUgras{mai sayar da
ƙwarƙwaran teku} \textsuperscript{(97)} wadda take sayar da
\VUgras{ƙwarƙwaran teku} \textsuperscript{(98)}. Tana buɗe su idan mutum
yana so.

%%K t15-16-1 p
16. --- Dukan waɗannan ƴan kasuwa suna biyan lasisi ko harajin wuri ;
shi ya sa \VUgras{ɗan sanda} \textsuperscript{(100)} yake korar
\VUgras{matan sayar da ganyayyaki masu yawo} \textsuperscript{(99)} masu
tausayi ba tare da tausayi ba, waɗanda suke yi musu takara suna kai
\VUgras{karas, da ƙananan raɗis, da turnif, da albasa mai tsawo} ko
\VUgras{ɗamin asfaragas, da ƙananan wake, da alayyaho, da yakuwa, da
ganyen ruwa} da \VUgras{ganyen ƙamshi} a kan ƙananan abin hawansu.

%%K t15-tit-8 sub
\VUsaut{4.0mm}
\VUpk{10.2pt}{40}{Kula da ɗan sanda !}
\VUsaut{3.0mm}

%%K t15-17-1 p
17. --- Yaron da yake sayar da \VUgras{tafarnuwa} \textsuperscript{(101)} da
\VUgras{albasa,} ya san sarai cewa shi ma, ba shi da izinin sayar da
kayansa kusa da kasuwar. Yana gudu, yana gaggawa cikin haɗarin kife
kwandon \VUgras{ƙaguwa,} da \VUgras{ƙaguwar ruwa} da \VUgras{jatan
lande} na \VUgras{mai sayar da} \textsuperscript{(102)} \VUgras{manyan jatan
lande} da \VUgras{jatan lande masu ƙafa.}

%%K t15-18-1 p
18. --- Kowa yana hargitsi yana kuma hayaniya ƙwarai ; amma wannan ba
ya hana a ji muryar \VUgras{mai sa gilashi} \textsuperscript{(103)} mai
yawo sarai, ko kiran \VUgras{mai sayar da gawayi} \textsuperscript{(104)}
wanda ya bar karusarsa na ɗan lokaci domin ya kai \VUgras{ƙaramin buhun
gawayi} \textsuperscript{(105)}.

\VUsaut{6.0mm}
\VUfilet{22mm}
